\documentclass[a4paper,12pt]{article}
\usepackage[utf8]{inputenc}
\usepackage[T1]{fontenc}
\usepackage[french]{babel}
\usepackage{amsmath,amssymb}
\usepackage{geometry}

\geometry{margin=2.5cm}
\title{R\'esolution dans $\mathbb{C}$ d'une \'equation irrationnelle}
\author{S.\ Jaubert}
\date{F\'evrier 2026}

\begin{document}
\maketitle

\section*{\'Enonc\'e}

R\'esoudre dans $\mathbb{C}$ :
\[
  \frac{1}{1+\sqrt{x^2-x+1}} + \frac{x}{1+\sqrt{x^2+x+1}} = 1
\]

\section*{Solution}

On pose $u = \sqrt{x^2-x+1}$ et $v = \sqrt{x^2+x+1}$.

\bigskip

\textbf{\'Etape 1 -- R\'eduction alg\'ebrique.}

En r\'eduisant au m\^eme d\'enominateur, l'\'equation devient :
\[
  (1+v) + x(1+u) = (1+u)(1+v)
\]
Apr\`es d\'eveloppement et simplification :
\[
  x(1+u) = u(1+v) \qquad (\star)
\]

\bigskip

\textbf{\'Etape 2 -- \'Elimination des radicaux.}

De $(\star)$ on tire $uv = x + u(x-1)$, soit $v = \dfrac{x}{u} + (x-1)$. En \'elevant au carr\'e et en utilisant $u^2 = x^2-x+1$ et $v^2 = x^2+x+1$, apr\`es calcul on obtient (pour $x \neq 0$, $x \neq 1$) :
\[
  3x^2 - 4x + 3 = 2(x-1)\,u
\]
En \'elevant \`a nouveau au carr\'e, on aboutit au polyn\^ome :
\[
  5x^4 - 12x^3 + 18x^2 - 12x + 5 = 0
\]

\bigskip

\textbf{\'Etape 3 -- Polyn\^ome palindromique.}

Les coefficients $(5,\, -12,\, 18,\, -12,\, 5)$ sont sym\'etriques. En divisant par $x^2$ et en posant $t = x + \dfrac{1}{x}$ :
\[
  5\!\left(t^2 - 2\right) - 12t + 18 = 0 \quad \Longrightarrow \quad 5t^2 - 12t + 8 = 0
\]
Le discriminant vaut $\Delta_t = 144 - 160 = -16 < 0$, d'o\`u :
\[
  t = \frac{6 \pm 2i}{5}
\]

\textbf{Cons\'equence :} l'\'equation n'admet \textbf{aucune solution r\'eelle}.

\bigskip

\textbf{\'Etape 4 -- R\'esolution des \'equations quadratiques.}

Pour chaque valeur de $t$, on r\'esout $x^2 - tx + 1 = 0$.

\medskip

\noindent\textit{Cas $t = \dfrac{6+2i}{5}$ :}
L'\'equation $5x^2 - (6+2i)\,x + 5 = 0$ a pour discriminant :
\[
  \Delta = (6+2i)^2 - 100 = -68 + 24i
\]
On v\'erifie que $\sqrt{\Delta} = \alpha + i\beta$, o\`u
\[
  \alpha = \sqrt{10\sqrt{13} - 34}, \qquad \beta = \sqrt{10\sqrt{13} + 34}, \qquad \alpha\,\beta = 12.
\]
Les deux racines sont :
\[
  x = \frac{6 + 2i \pm (\alpha + i\beta)}{10}
\]

\medskip

\noindent\textit{Cas $t = \dfrac{6-2i}{5}$ :} par conjugaison, les racines sont $\overline{x}$.

\bigskip

\textbf{\'Etape 5 -- \'Elimination des solutions parasites.}

La mise au carr\'e a pu introduire des racines \'etrang\`eres. La v\'erification num\'erique montre que seules les racines correspondant au signe $+$ dans chaque cas satisfont l'\'equation originale. On obtient le couple de solutions conjugu\'ees :

\[
\boxed{x = \frac{6 + \alpha}{10} \pm \, i\;\frac{2 + \beta}{10}}
\]

o\`u $\alpha = \sqrt{10\sqrt{13} - 34}$ et $\beta = \sqrt{10\sqrt{13} + 34}$.

\bigskip

Num\'eriquement : $\;x \approx 0{,}743 \pm 1{,}037\,i$.

\end{document}
